\chapter{Główne twierdzenia}

W tym rozdziale przedstawiamy kluczowe wyniki teoretyczne dotyczące markowskich procesów decyzyjnych z quasi-hiperbolicznym dyskontowaniem. Twierdzenia te stanowią fundamenty teoretyczne dla algorytmów przedstawionych w kolejnych rozdziałach.

\section{Twierdzenie o redukcji polityk}

Pierwszym kluczowym wynikiem jest pokazanie, że złożone polityki zależne od historii można zredukować do prostszej formy.

\begin{theorem}[Redukcja polityk]
\label{thm:policy-reduction}
Dla dowolnej $\phi \in \Phi$ istnieją reguły $\mu$ i $\phi_s$ takie, że
\[
  J_{\phi}^{\alpha,\beta}(s) = J_{(\mu,\phi_s)}^{\alpha,\beta}(s), \qquad \forall s \in S.
\]
\end{theorem}

Pierwszy krok steruje dowolna reguła $\mu = \phi_0$, kolejne kroki wykorzystują stacjonarną $\phi_s$. Analiza optymalności sprowadza się do klas polityk dwuskładnikowych $(\mu,\phi_s)$.

\begin{proof}
Dla dowolnej historycznej polityki $\phi = (\phi_0, \phi_1, \ldots)$ mamy
\[
  J_{\phi}^{\alpha,\beta}(s) = \mathbb{E}\!\left[ \sum_{t=0}^{\infty} d(t)\, r(s_t, a_t) \;\middle|\; s_0 = s \right],
\]
gdzie $d(0) = 1$ oraz $d(t) = \alpha\,\beta^{t}$ dla $t \ge 1$.

Rozwijając wartość oczekiwaną otrzymujemy
\[
  J_{\phi}^{\alpha,\beta}(s) = \mathbb{E}\!\left[ r(s_0,a_0) + \sum_{t=1}^{\infty} \alpha \beta^{t}\, r(s_t,a_t) \;\middle|\; s_0 = s \right].
\]

Warunkując po pierwszej akcji $a_0 \sim \phi_0(\cdot \mid s)$ dostajemy
\[
  J_{\phi}^{\alpha,\beta}(s) = \mathbb{E}_{a_0 \sim \phi_0(\cdot \mid s)}\!\left[ r(s,a_0) + \,\mathbb{E}\!\left[ \sum_{t=1}^{\infty} \alpha \beta^{t}\, r(s_t,a_t) \;\middle|\; s_0 = s, a_0 \right] \right].
\]

Kolejne warunkowanie względem $s' \sim q(\cdot \mid s,a_0)$ daje
\[
  J_{\phi}^{\alpha,\beta}(s) = \mathbb{E}_{\substack{a_0 \sim \phi_0(\cdot \mid s) \\ s' \sim q(\cdot \mid s,a_0)}}\!\left[ r(s,a_0) + \alpha\beta\,\mathbb{E}\!\left[ \sum_{t=1}^{\infty} \beta^{t-1}\, r(s_t,a_t) \;\middle|\; s_1=s' \right] \right].
\]

Pozostała suma dyskontowana ma postać wykładniczą, więc
\[
  \mathbb{E}\!\left[ \sum_{t=1}^{\infty} \beta^{t-1}\, r(s_t,a_t) \;\middle|\; s_1 = s' \right] = \,J_{\overrightarrow{\phi}}^{\beta}(s_1),
\]
gdzie $\overrightarrow{\phi} = (\phi_1, \phi_2, \ldots)$ i $J^{\beta}$ odpowiada klasycznemu dyskontowaniu wykładniczemu.

Dla takiego dyskontowania istnieje stacjonarna polityka $\phi_s \in \Phi_s$, która spełnia
\[
  J^{\beta}_{\overrightarrow{\phi}}(s') = J^{\beta}_{\phi_s}(s'), \qquad \forall s' \in S.
\]

Podstawiając do poprzedniego równania, otrzymujemy
\[
  J_{\phi}^{\alpha,\beta}(s) = \sum_{a}\phi_{0}(a \mid s)\!\left[ r(s,a) + \alpha\,\beta\, \sum_{s'} q(s' \mid s,a) J^{\beta}_{\phi_s}(s') \right].
\]

Wybierając $\mu = \phi_0$ jako regułę pierwszego kroku, uzyskujemy żądaną równość
\[
  J_{\phi}^{\alpha,\beta}(s) = J_{(\mu,\phi_s)}^{\alpha,\beta}(s), 
\]
co kończy dowód.
\end{proof}

\section{Lemat o kontrakcji}

Następny wynik pokazuje, że operator Bellmana dla polityki stacjonarnej jest kontrakcją.

\begin{theorem}[Kontrakcja operatora Bellmana dla QH]
\label{thm:bellman-contraction}
Definiujemy operator $T^{\phi_s}: \mathbb{R}^{|S|} \to \mathbb{R}^{|S|}$:
\begin{equation}
(T^{\phi_s} J)(s) = \sum_a \phi_s(a\mid s) \left[ r(s, a) + \sum_{s'} q(s'\mid s, a) \left[ -(1 - \alpha)\beta \sum_{a'} \phi_s(a'\mid s') r(s', a') + \beta J(s') \right] \right].
\end{equation}
Operator $T^{\phi_s}$ jest kontrakcją w normie $\max$ z faktorem $\beta < 1$. Jedyny punkt stały $J = T^{\phi_s} J$ to $J^{\alpha, \beta}_{\phi_s}$.
\end{theorem}

\begin{proof}
Dla dowolnych $J_1, J_2 \in \mathbb{R}^{|S|}$ mamy:
\[
\|T^{\phi_s} J_1 - T^{\phi_s} J_2\|_\infty \leq \beta \|J_1 - J_2\|_\infty,
\]
co wynika wprost z twierdzenia Banacha o odwzorowaniu zwężającym. Ponieważ $\beta < 1$, operator jest kontrakcją, a zatem posiada dokładnie jeden punkt stały.
\end{proof}

\section{Twierdzenie o zbieżności polityk}

\begin{theorem}[Zbieżność iteracji polityk]
\label{thm:policy-convergence}
Niech $(J_n)$ będzie ciągiem zdefiniowanym rekurencyjnie przez
\[
J_{n+1} = T^{\phi_s} J_n, \quad n = 0, 1, 2, \ldots
\]
dla dowolnego $J_0 \in \mathbb{R}^{|S|}$. Wówczas
\[
\lim_{n \to \infty} J_n = J^{\alpha, \beta}_{\phi_s}.
\]
\end{theorem}

\begin{proof}
Z Twierdzenia~\ref{thm:bellman-contraction} wiemy, że $T^{\phi_s}$ jest kontrakcją z faktorem $\beta < 1$. Z twierdzenia Banacha o punkcie stałym dla przestrzeni zupełnych wynika, że dla dowolnego punktu startowego $J_0$, iteracje $J_{n+1} = T^{\phi_s} J_n$ zbiegają do jedynego punktu stałego $J^{\alpha, \beta}_{\phi_s}$.

Dokładniej, mamy oszacowanie błędu:
\[
\|J_n - J^{\alpha, \beta}_{\phi_s}\|_\infty \leq \beta^n \|J_0 - J^{\alpha, \beta}_{\phi_s}\|_\infty,
\]
co pokazuje wykładniczą szybkość zbieżności.
\end{proof}

\section{Twierdzenie o istnieniu optymalnej polityki deterministycznej}

\begin{theorem}[Deterministyczna optymalna polityka]
\label{thm:deterministic-optimal}
Istnieje deterministyczna optymalna para $(\mu^\star, \phi_s^\star)$ taka, że
\begin{equation}
V^{\alpha, \beta}_{\mu^\star, \phi_s^\star}(s) = \max_{\mu \in \Phi_m} \max_{\phi_s \in \Phi_s} V^{\alpha, \beta}_{\mu, \phi_s}(s), \quad \forall s \in S.
\end{equation}
\end{theorem}

\begin{proof}
Z Twierdzenia~\ref{thm:policy-reduction} wiemy, że wystarczy rozważać polityki postaci $(\mu, \phi_s)$.

Punktem wyjścia jest redukcja polityk:
\begin{align}
V^{\alpha, \beta}_{\mu^\star, \phi_s^\star}(s) &= \max_{\mu \in \Phi_m} \max_{\phi_s \in \Phi_s} \sum_{a \in A} \mu(a\mid s) \Big[ r(s, a) + \alpha \beta \sum_{s' \in S} q(s'\mid s, a) V^{\beta}_{\phi_s}(s') \Big] \nonumber \\
&= \max_{\mu \in \Phi_m} \sum_{a \in A} \mu(a\mid s) \Big[ r(s, a) + \alpha \beta \sum_{s'} q(s'\mid s, a) \max_{\phi_s \in \Phi_s} V^{\beta}_{\phi_s}(s') \Big].
\end{align}

Dla dyskontowania wykładniczego optymalna jest $\pi_\beta^\star$, stąd $\phi_s^\star = \pi_\beta^\star$ oraz $V^{\beta}_{\phi_s^\star} = V^{\beta}_\star$.

Optymalna reguła pierwszego kroku:
\begin{equation}
\mu^\star(s) = \arg\max_{a \in A} \left[ r(s, a) + \alpha \beta \sum_{s'} q(s'\mid s, a) V^{\beta}_{\pi^\star_\beta}(s') \right].
\end{equation}

Polityka $\phi_s^\star$ jest deterministyczna, ponieważ $\pi_\beta^\star$ w klasycznym MDP taką jest. Optymalizacja $\mu$ nad sympleksem $\Pr(A)$ jest liniowa, więc optimum znajduje się w wierzchołku, co oznacza deterministyczne $\mu^\star$.

W praktyce wystarczy przeszukiwać $\mu^\star \in A^S$ oraz $\phi_s^\star \in A^S$.
\end{proof}

\section{Twierdzenie o zbieżności QH Q-Learning}

Ostatnie twierdzenie zapewnia zbieżność algorytmu uczenia się bez modelu.

\begin{theorem}[Konwergencja QH Q-Learning]
\label{thm:qh-qlearning-conv}
Załóżmy, że kroki uczenia $(\eta_n)$ i $(\theta_n)$ spełniają warunki Robbins--Monro:
\[
\sum_n \eta_n = \infty, \quad \sum_n \eta_n^2 < \infty, \quad \sum_n \theta_n = \infty, \quad \sum_n \theta_n^2 < \infty,
\]
oraz dodatkowo $\theta_n / \eta_n \to 0$ (dwie skale czasowe). 

Wtedy ciągi $(Q_n, W_n)$ generowane przez Algorytm QH Q-Learning zbiegają prawie na pewno do $(Q^{\alpha, \beta}_{\mu^\star, \phi_s^\star}, Q^{\beta}_{\phi_s^\star})$. 

Polityki ekstrahowane jako $\mu^\star(s) = \arg\max_a Q_n(s,a)$ oraz $\phi_s^\star(s) = \arg\max_a W_n(s,a)$ są optymalne.
\end{theorem}

\begin{proof}[Szkic dowodu]
Dowód wykorzystuje teorię aproksymacji stochastycznej o dwóch skalach czasowych:

\textbf{Szybka skala ($W_n$):} 
Aktualizacja $W_n$ to klasyczny Q-learning dla dyskontowania wykładniczego:
\[
W_{n+1}(s, a) \leftarrow W_n(s, a) + \eta_n \left[ r(s, a) + \beta \max_{a'} W_n(s', a') - W_n(s, a) \right].
\]
Z klasycznej teorii Q-learningu wynika, że $W_n \to Q^{\beta}_{\phi_s^\star}$ prawie na pewno.

\textbf{Wolna skala ($Q_n$):}
Aktualizacja $Q_n$ ma postać:
\[
Q_{n+1}(s, a) \leftarrow Q_n(s, a) + \theta_n \left[ (1 - \alpha) r(s, a) + \alpha W_n(s, a) - Q_n(s, a) \right].
\]
Ponieważ $\theta_n / \eta_n \to 0$, proces $W_n$ działa w szybszej skali czasowej i można go traktować jako quasi-stacjonarny z punktu widzenia $Q_n$.

W granicy $W_n \approx Q^{\beta}_{\phi_s^\star}$, więc aktualizacja $Q_n$ zbieżna jest do układu liniowego ze stałym punktem $Q^{\alpha, \beta}_{\mu^\star, \phi_s^\star}$.

\textbf{Zachowanie greedy:}
Polityki greedy względem zbieżnych funkcji Q są optymalne:
\begin{align}
\mu^\star(s) &= \arg\max_a Q^{\alpha, \beta}_{\mu^\star, \phi_s^\star}(s,a), \\
\phi_s^\star(s) &= \arg\max_a Q^{\beta}_{\phi_s^\star}(s,a).
\end{align}

Formalne uzasadnienie wymaga zastosowania teorii ODE (ordinary differential equations) dla aproksymacji stochastycznej oraz analizy dwuskalowej, co wykracza poza zakres niniejszej pracy. Szczegóły można znaleźć w pracy \cite{eshwar2024}.
\end{proof}
